% ================================================================
%  chapters/parts/ch06-b.tex
%  第 6 章  催化过程的拉曼光谱方法（分片 b）
%  源：S1（上册）第 6 章，书内 p.373—412 / PDF p.392—431
%  覆盖：§6.4 紫外拉曼光谱仪在催化材料中的应用研究（含 6.4.1—6.4.6）
%        与 §6.5 展望；另含 §6.3.6 末段（原位反应研究实例）
%  分片 a 已写 §6.1—§6.2（原理、实验技术、荧光发生机制），本片不重复。
% ================================================================

\subsection{拉曼光谱在催化剂表征中的应用}

6.1—6.2 节（见 \chapref{ch:06}）已经说明了三件事：拉曼测的是分子极化率变化，荧光为什
么能把拉曼信号整个淹没，以及一台光谱仪的五个部件各自把什么条件写进了谱图。从本节起只问
一个问题：\textbf{面对一块真实的催化剂，谱图上能读出哪几个量，这些读数在什么条件下失效}。

可读的独立量只有三个——\textbf{峰位、峰宽与相对强度}；本章还额外提供了一个维度：\textbf{信
号来自多深}（紫外拉曼靠吸收深度把表面与体相分开，共焦拉曼靠针孔把不同深度分开，见后两节）。
拉曼之所以在催化里站得住，靠的是两条与红外、X 射线衍射都不同的性质：载体本身（\ce{SiO2}、
$\gamma$-\ce{Al2O3}）的拉曼散射截面很小，负载物种的谱峰不会被载体淹没；水的拉曼散射截面
非常小，因此水相与水热体系可以直接测、而且液相信号与固相信号能同时拿到
\src{S1 §6.4.6.3，书内 p.388 / PDF p.407}。

\begin{probout}
\textbf{输入}：一块处于给定预处理历史与给定环境（温度、气氛、介质）下的催化剂样品，加上一
条已知波长与功率密度的激发线。\textbf{输出}：散射光强度对拉曼位移的曲线，以及（若用共焦或
紫外激发）这条曲线对应的采样深度。

\textbf{本节要回答的问题}是：这条曲线上有哪些信息是\textbf{可以}读的，哪些是\textbf{不能}
读的。可以读的是——骨架振动与物相（四元环、双六元环呼吸振动；\ce{MoO3}、\ce{ZrO2}、
\ce{TiO2} 各晶相的特征峰）、负载物种的配位环境（端基 Mo—O 与桥键 Mo—O—Mo）、吸附态
与孔道内分子的扰动（C—H 振动区的位移）、以及随温度或时间发生的相对强度变化。不能读的是
——绝对浓度（强度只可相对比较，且共振与增强效应会按模式改变比例）、"峰不存在所以物种不存
在"（脱水、覆盖、吸收深度都会让峰消失）、以及"谱峰位置唯一对应某个物种"（多物种在同一区
间都有峰，见下文）。
\end{probout}

\subsubsection{骨架、负载物种与吸附态：三类可读信息}

\textbf{（1）骨架与物相。}分子筛的骨架振动落在 $200$—\qty{1200}{cm^{-1}} 的低频区，四元环
与双六元环的呼吸振动分别位于约 \qty{500}{cm^{-1}} 与 \qty{298}{cm^{-1}}、\qty{380}{cm^{-1}}
附近：无定形前驱体只给出宽包，晶化后变成可分辨的锐峰（原位实例见后面紫外拉曼一节）。氧化物
给出更硬的结构判据——\ce{MoO3} 晶相在约 \qty{820}{cm^{-1}} 有特征峰，它的\textbf{缺失}才是
"钼以高分散表面物种存在、没有晶相 \ce{MoO3} 析出"的证据。

这里有一条贯穿全章的判据纪律：\textbf{谱峰位置只能给出"存在某种配位环境"，不能唯一确定物
种}。源文本在 Mo、V、W 混合氧化物催化剂的共焦拉曼研究中把这一点说得很直白：在该区域内
Mo、V 和 W 都有拉曼谱峰，峰的位置取决于它们的配位环境和对称性，因而\textbf{将这些谱峰直接
归属于某一特定物种是不可能的}\src{S1 §6.4.5，书内 p.381 / PDF p.400}。要从"某个波数有峰"
走到"某个物种存在"，必须另加约束：换激发波长看共振选择性（下一节）、用纯晶相参考谱比对、或
对成百上千条谱做多变量解析。

\textbf{（2）负载物种与载体相互作用。}低负载量是拉曼最能发挥作用的区间：载体几乎不散射，
因而表面物种的弱峰不会被人为放大或掩盖。源文本指出，由于可见拉曼光谱的灵敏度低和荧光干扰，
早期研究主要集中在较高负载量催化剂上，而\textbf{只有研究低负载量催化剂才能够给出表面物种与
载体相互作用的直接信息}\src{S1 §6.4.6.2，书内 p.385 / PDF p.404}。

\textbf{（3）吸附态与孔道内的分子。}两个方向的例子。正向的例子：二甲苯吸附于 ZSM-5 分子
筛后，与纯二甲苯的主要不同在 C—H 振动范围内，吸附后这些谱峰向高频方向位移，作者归因于
C—H 振动模式\textbf{受到分子筛骨架的限制}\src{S1 §6.4.3，书内 p.377 / PDF p.396}。反向
的例子更重要——\ce{AlPO4}-18 和 SAPO-34 分子筛经焙烧脱水后在拉曼谱图中难以观察到其骨架振动谱
峰，吸附微量水后骨架振动谱峰又出现了\src{S1 §6.4.3，书内 p.377 / PDF p.396}："谱图上没有峰"同样是一条
需要解释的信息，它可能来自样品状态（含水量、表面覆盖、物相）而不是结构的不存在。因此任何
"某物种消失"的结论都必须连同样品的预处理与含水量一起报告。

\emph{核对说明}：已对照原页面图像（书内 p.377 / PDF p.396）。原书此处举的是\textbf{两个}例子——
「Huang$^{[146]}$ 用 FT-Raman 研究了二甲苯在 ZSM-5 分子筛上的吸附……吸附后这些部分的谱峰向高频
方向位移，这是由 C—H 振动模式受到分子筛骨架限制而引起的」，以及「Löffler 和 Bergmann$^{[150]}$
研究了焙烧和脱水对分子筛骨架的影响，他们发现 \ce{AlPO4}-18 和 SAPO-34 分子筛经焙烧脱水后在拉曼
谱图中难以观察到其骨架振动谱峰，吸附微量水后，如图 6-20 所示其骨架振动谱峰又出现了」。图 6-20
题注为「\ce{AlPO4}-18（a、b）和 SAPO-34（c、d）的 FT-Raman 光谱$^{[150]}$」，a、c 为焙烧脱水，
b、d 为焙烧脱水后再部分吸水。作者姓名原书印作 \textbf{Löffler 和 Bergmann}（源 OCR 作「L6mer
和 Bergmalm」，系 OCR 失真）。原书\textbf{只报告现象并将其归属为骨架振动，未给出该现象的机制}；
本分片"谱图上没有峰同样是一条需要解释的信息"及其后的报告要求为本分片补写的判断，非原书文字。
另：源 OCR 文本层存在「钼/铝」互讹（如 OCR 作「$\gamma$-\ce{Al2O3} 负载的氧化铝催化剂」），
本分片按化学式与上下文取「钼」。

\subsubsection{原位拉曼与原位池：温度、气氛与窗口材料}
\label{sec:ch06b-insitu}

原位（in situ）是拉曼相对其他表征手段最大的优势之一：它\textbf{较容易实现原位检测，可在高
温高压下表征催化剂的变化，且不受气相或载体的干扰}\src{S1 §6.5，书内 p.402 / PDF p.421}。
把样品从反应环境中取出来再测，对高温高压水热体系尤其危险——在反应中将产物取出可能导致中间
物种的结构性质发生明显变化，失去研究价值\src{S1 §6.4.6.3，书内 p.388 / PDF p.407}。

源文本在本节开头承接上一节的“原位反应研究”，给出一个原位激光拉曼的实例：Lunsford 等\src{S1 §6.3.6 末段，书内 p.373 / PDF p.392}用原位激光拉曼光谱详细研究了 \ce{BaO}/\ce{MgO} 催化剂上 \ce{NO} 和 \ce{N2O} 的降解反应，给出的 \ce{N2O} 降解反应步骤为：① \ce{N2O} 和表面 \ce{O^{2-}} 反应给出 \ce{N2} 和表面 \ce{O2^{2-}} 物种；② 两个表面 \ce{O2^{2-}} 物种反应给出 \ce{O2} 和两个表面 \ce{O^{2-}}；③ \ce{O2} 和两个表面 \ce{O^{2-}} 反应给出非活性的 \ce{O2^{-}} 物种。

\emph{核对说明}：已对照原页面图像（书内 p.373 / PDF p.392），本节三步步骤与原书逐字一致。源 OCR 把 \ce{BaO}/\ce{MgO} 识作「ＢａＮＺｏ」、把 \ce{N2O} 识作「ＮＺｏ」、把三种表面氧物种的电荷符号识作「０２一」「０１一」，均在页面图像上复原；原书亦未给出更多机理细节，本分片照录，不代为补写。该实例位于 §6.3.6 末段（书内 p.373），前文原引 §6.4.6.3 有误，此处按其真实出处标注。

\textbf{原位池的三个设计变量。}

\textbf{温度与压力。}源文本描述的水热合成原位拉曼池被铜制加热线圈环绕，能够耐受
\qty{250}{\degC} 的高温与 \qty{40}{bar} 的高压（$1$ bar $=10^5$ Pa）\src{S1 §6.4.6.3，
书内 p.388 / PDF p.407}。相比之下，近红外傅里叶变换拉曼的可用温度低得多：源文本明确指出近
红外波段拉曼光谱仪在原位催化研究中受到限制，\textbf{因为样品不能升太高温度（$<\qty{150}{\degC}$）}，
而且许多实用催化剂在近红外区也有荧光\src{S1 §6.4.3，书内 p.378 / PDF p.397}。同一台样品池
在紫外、可见、近红外三条路线上的温度上限并不相同，这是一条容易忽略的选型约束。

\textbf{窗口材料与密封。}通光部分通常选熔融石英（紫外到近红外透明、杂质少、紫外下不发光，
见本章仪器的部件表）。水热体系里更关键的是\textbf{窗口的形状}：源文本的做法是\textbf{用一个
透镜作窗口}——透镜通过硅橡胶密封于原位池的顶部，既当窗口又当聚焦元件，这样获得的拉曼信号
是\textbf{用平面镜作窗口时的 3—4 倍}\src{S1 §6.4.6.3，书内 p.388 / PDF p.407}。该池另配聚四
氟乙烯池与聚四氟乙烯内衬（源文本图注），用于耐腐蚀的水热碱性介质。表~\ref{tab:insitu-window}
把这几条方案与它们各自的代价并排放。

两端都要看：透镜窗口换来 3—4 倍信号，但它把有机密封件的耐温上限带进了池体指标；石英在高温
水热碱性介质中会被腐蚀，这正是需要聚四氟乙烯内衬的原因。选窗口不是选"哪个透明"，而是选"哪个
能在实验条件下既透明又不参与反应"。

\textbf{气氛。}惰性、氧化、还原、水蒸气四种气氛对应四类完全不同的读数：在还原气氛（\ce{H2}）
中，表面钼、钒、铁的价态与配位会随温度改变，谱峰变化反映的是\textbf{气氛与温度写进样品的
状态}，不是"结构自己变了"。因此原位谱的图注必须写清气氛组成、流量与升降温程序，否则同一峰
位移无法在两组实验之间比较。

\begin{table}[htbp]
\centering
\caption{原位拉曼池的通光与耐环境方案（本节实例）。读者应看"代价"一列：源文本给出的每一种
方案都是用一项能力换另一项能力，没有一种窗口同时满足紫外透明、耐腐蚀、耐高温高压与高收集效
率。\src{据 S1 §6.4.6.3，书内 p.388—389 / PDF p.407—408 归纳。}}
\label{tab:insitu-window}
{\srcsize
\begin{tabularx}{\columnwidth}{@{}lXX@{}}
\toprule
方案 & 源文本给出的做法与数据 & 代价与限制 \\
\midrule
平面镜作窗口 & 可获得拉曼信号，作为该池的参照方案 & 信号只有透镜窗口方案的 $1/3$—$1/4$；无聚焦能力，照射面积大、功率密度低 \\
透镜作窗口（硅橡胶密封于池顶） & 既作窗口又作聚焦元件，把激光聚焦到被研究样品表面；信号为平面镜方案的 3—4 倍；池体可耐 \qty{250}{\degC} 与 \qty{40}{bar} & 有机密封件是有机物，会贡献荧光背底并限制温度上限；不耐强腐蚀介质；换波长或换焦距时透镜位置需重新对准 \\
聚四氟乙烯池与内衬 & 用于水热（碱性）介质中的耐腐蚀 & 聚四氟乙烯在高温下会软化、透气，机械强度低，不能单独承担高压 \\
熔融石英通光件 & 紫外到近红外全波段透明、杂质少、紫外下不发光（见本章仪器部件表） & 高温水热碱性介质下会被腐蚀；紫外下长期照射可能出现色心而增大吸收 \\
\bottomrule
\end{tabularx}}
\end{table}

\emph{核对说明}：已对照原页面图像（书内 p.388—389 / PDF p.407—408）。原书
§6.4.6.3 记：原位池被铜制加热线圈环绕，一个透镜通过硅橡胶密封于池顶，用来把激光聚焦于样品表面；
「利用透镜作为原位池的窗口所获得的拉曼信号是利用平面镜作为窗口获得拉曼信号的 3$\sim$4 倍」；
在分子筛合成的水热条件下可耐 \qty{250}{\degC} 与 \qty{40}{bar}（1 bar
$=10^{5}$ Pa）。本表「平面镜作窗口」行的 1/3—1/4 即该 3$\sim$4 倍的倒数表述，
「透镜作窗口」行的聚焦作用、信号倍数与耐温耐压值与原书一致。图 6-29 的标注给出透镜窗口、
聚四氟乙烯池与聚四氟乙烯线反应釜；深/浅两种样品池（分别聚焦液相与液固界面）亦见该图。
聚四氟乙烯与熔融石英两行的材料性质，以及各行「代价与限制」中的有机密封件荧光背底、聚四氟乙烯高温软化透气、
石英紫外色心等，属通用材料学常识，原书未逐条给出。


\begin{measure}
\textbf{原位拉曼必须记下来的六个条件}（缺任何一条，谱图就无法与别人的谱图比较）：
\begin{itemize}
\item \textbf{激发波长与功率密度}：不只是激光器的标称功率，而是\textbf{样品上的功率密度}
  （功率除以光斑面积）。功率密度决定样品被加热的程度，是原位实验最容易失控的一项。
\item \textbf{光束尺寸与照射方式}：静止照射、旋转样品池、线扫描或散焦，四者的局部温升差别
  很大。本章仪器部分已指出，可旋转样品池是避免样品被打坏的常用对策。
\item \textbf{采样时间与是否累加}：光致变化是时间的函数，只给总积分时间而不给单次曝光时间，
  无法判断样品是否在采集过程中被改变。
\item \textbf{池体温度与压力及其测量点}：测温点在与不在被照射区，结论完全不同；本章关于
  其他原位技术的经验（测温误差直接进入构效关系）在这里同样适用。
\item \textbf{气氛组成与流量}：还原性气氛下价态与配位随温度变化，谱峰位移不能跨气氛比较。
\item \textbf{窗口/密封材料与样品预处理历史}：聚合物密封件贡献荧光背底；含水量与焙烧历史
  直接决定骨架峰是否出现（见上一小节 SAPO-34 的例子）。
\end{itemize}
把六条一起报出来，原位谱才是证据；只报"在某温度下测得的原位拉曼谱"，它只是图。
\end{measure}

\textbf{激光功率导致的样品加热与光致变化}是原位拉曼最需要主动排除的干扰。源文本把这件事直
接列为本征拉曼的限制：由于\textbf{激光导致的样品脱水、相变、部分还原，甚至完全降解}，以及
荧光干扰和固有的低灵敏度，仍然限制了拉曼光谱在原位表征中的进一步应用
\src{S1 §6.4.1 末，书内 p.374 / PDF p.393}。链条是：激光功率密度 $\to$ 局部温升（深色样品
在可见区吸收强，升温尤其快）$\to$ 脱水、相变、还原、降解 $\to$ 谱峰位移与强度改变。麻烦在
于\textbf{这条链条的产物与"真实结构变化"在谱图上长得一模一样}，唯一的区分办法是主动改变激
光功率重测。

\begin{pitfall}
\textbf{把激光加热导致的谱带位移当成结构变化}，是原位拉曼最常见、也最难自查的一类误判。三
条判据，缺一条就不能下结构结论：
\begin{itemize}
\item \textbf{功率依赖}：把激光功率降到原来的一半、四分之一各测一次。若峰位随功率连续回移、
  峰宽随功率变窄，说明位移来自热效应；若峰位不随功率改变，只随温度或气氛改变，才是样品状态
  的变化。
\item \textbf{时间依赖}：在固定功率下连续采集同一位置。若峰位在几十秒内单调漂移，说明样品正
  在被逐步改变（脱水、积炭、还原），此时"随时间演化的原位谱"记录的是激光的破坏过程。
\item \textbf{可逆性}：撤去激光、回到同一温度后重测。可逆的是热效应；不可逆的位移意味着发生
  了光致化学变化（脱水、部分还原、降解），该样品不能再用这条功率密度测量。
\end{itemize}
特别要注意三类高危样品：深色样品（可见区强吸收）、含水或含羟基样品（脱水会改变骨架与表面
物种）、以及高比表面氧化物（光致还原在紫外激发下更易发生）。三者都是催化研究的常客，因此这
条误判不是边缘情形，而是默认要排除的第一项。
\end{pitfall}

\subsubsection{近红外路线：傅里叶变换拉曼的收益与代价}

解决荧光干扰有两条"把激发波长移开"的路线：向近红外（$>\qty{700}{nm}$）移动可以避开大部分
荧光，因为近红外光频率低，一般不会激发电子吸收带；向深紫外（$<\qty{300}{nm}$）移动则靠能
级间隙把拉曼与荧光在波长上分开（下一节）。近红外路线的代表是傅里叶变换拉曼光谱
（FT-Raman）。

FT-Raman 一般采用波长 \qty{1064}{nm} 的近红外激发光源，分子的荧光不被激发，因此可以避免
许多样品拉曼光谱中的荧光干扰；仪器以近红外激光为激发光源并引进傅里叶变换红外光谱仪中常用的
傅里叶变换技术，主要由光源、迈克耳孙干涉仪和检测器组成，其中干涉仪是最重要的部件，把分束器
换成石英型以利于近红外光透过\src{S1 §6.4.3，书内 p.377 / PDF p.396}。

它的代价源文本列了四条，逐条都能量化：
\begin{itemize}
\item \textbf{①散射强度按频率四次方下降。}散射光强度与激发光频率的四次方成正比，采用频率较
  低的红外光会使检测灵敏度大幅度降低。把激发线从 \qty{532}{nm} 换到 \qty{1064}{nm}，仅频率
  因子就使强度降到 $(532/1064)^4=1/16$。
\item \textbf{②探测器落后两个数量级。}红外区的探测器还很不成熟，与可见区的光电倍增管相比，
  灵敏度至少差两个数量级。①与②叠加，意味着近红外路线的总灵敏度损失远大于单看任一条。
\item \textbf{③样品热辐射本身成为干扰。}样品的热辐射对红外波段的拉曼光谱有很大的干扰——这
  一条直接把近红外路线与原位高温实验对立起来。
\item \textbf{④荧光并未被彻底避开。}一些稀土离子和过渡金属离子的电子态跃迁处于近红外区域，
  因此仍然无法避开荧光。
\end{itemize}
源文本的结论是：由于以上原因，近红外波段的拉曼光谱仪的应用受到限制，特别是在原位催化研究中
受到限制，因为样品不能升太高温度（$<\qty{150}{\degC}$），许多实用催化剂在近红外区也有荧光
\src{S1 §6.4.3，书内 p.378 / PDF p.397}。

一句话对比两条"移开波长"的路线：\textbf{近红外是用灵敏度换荧光，紫外是用选择性与探测深度换
荧光与灵敏度}。前者换得越来越不划算，后者才是本章剩余篇幅的主角。

\subsection{紫外拉曼与共振拉曼}

\subsubsection{共振拉曼效应：截面为什么会在吸收带内爆炸}
\label{sec:ch06b-resonance}

常规拉曼实验中使用的激发线波长远离化合物的电子吸收谱带，由此得到的光谱形状与激发波长的关系
只是服从激光频率的四次方关系。但当激发光源接近或落在化合物的电子吸收光谱带内时，某些拉曼谱
带的强度将大大增强，这种现象称为\textbf{共振拉曼效应}（resonance Raman effect），其理论基础
由 Kramers、Heisenberg 和 Dirac 建立\src{S1 §6.4.2，书内 p.375 / PDF p.394}。

\begin{principle}
\textbf{共振的分母。}共振拉曼效应来自跃迁极化率表达式中分母的趋零：当激光的频率 $\nu_0$ 与
吸收的频率 $\nu_{\mathrm{m}}$ 接近时，跃迁极化率的第一项分母趋于 $0$，跃迁极化率在形式上趋
于无穷大，因此共振拉曼谱带的强度将大大增强。可将拉曼散射截面的公式简化为

\begin{equation}
\sigma \;\propto\; \frac{1}{\left(h\nu_0-h\nu_{\mathrm{m}}\right)\left(h\nu_0+h\nu_{\mathrm{m}}\right)}
\label{eq:resonance-cross}
\end{equation}

式中 $\sigma$ 为拉曼散射截面，$h\nu_0$ 与 $h\nu_{\mathrm{m}}$ 是两个能量因子：\textbf{一个
是入射激光的光子能量，另一个是电子基态与激发态之间的能级差}。当两者相等时差项
$(h\nu_0-h\nu_{\mathrm{m}})$ 趋于零，截面急剧增大——这就是"共振"的全部内容。

\textbf{物理图像}（源文本描述，本分片未据原页面图像重绘）：当激发频率接近或重合于分子的一个
电子吸收带时，电子的跃迁过程由基态跃迁到激发态的\textbf{振动中间态}上；常规过程的\textbf{虚
拟}中间态变为\textbf{真实}的中间态，其跃迁概率极大提高，因而拉曼有效散射截面异常增大，某一
个或几个特定的拉曼带强度急剧增加，一般比非共振拉曼带强度增大 $10^{4}$—$10^{6}$ 倍，甚至出现正
常拉曼效应中观察不到的泛频及组合振动频率的谱峰。

\textbf{成立条件与代价}：式~\eqref{eq:resonance-cross} 是共振行为的\textbf{比例式}，不是截面
的完整表达式（完整的 Kramers—Heisenberg—Dirac 求和还含有阻尼项与各中间态的跃迁矩）。它只在
"激发线进入电子吸收带"时才有意义；一旦进入吸收带，被增强的只是\textbf{与该电子跃迁耦合的那
几个振动模式}，其余模式的增强因子各不相同。

\emph{核对说明}：已对照原页面图像。式~\eqref{eq:resonance-cross} 与书内 p.376 / PDF p.395 印作
$\sigma\propto 1/[(h\nu_0-h\nu_{\mathrm{m}})(h\nu_0+h\nu_{\mathrm{m}})]$ 逐字一致，书内
p.375 / PDF p.394 也确有「可将拉曼散射截面的公式简化为」一句引出该式。符号定义存在
\textbf{原书自身问题}：书内 p.375 / PDF p.394 写「当激光的频率 $\omega_0$ 与吸收的频率
$\omega_{ri}$ 接近时」（$\omega_0$ 为激光频率），而书内 p.376 / PDF p.395 的说明句写
「$h\nu_0$ 为电子基态和激发态的能级差，而 $h\nu_{\mathrm{m}}$ 为基态和虚态的能级差，即入射
激光的能量」——两页对同一对因子的赋值互换。本分片按 p.375 的定义书写（$\nu_0$ 为激光频率、
$\nu_{\mathrm{m}}$ 为电子跃迁频率），这也是通行写法。分母的完整形式确已印出：书内 p.375 /
PDF p.394 给出 Kramers—Heisenberg—Dirac 跃迁极化率
$(\tilde{\alpha}_{\rho\sigma})_{fi}=\frac{1}{h}\sum_{r\neq i,f}\left[\frac{\langle f|\hat{p}_{\rho}|r\rangle\langle r|\hat{p}_{\sigma}|i\rangle}{\omega_{ri}-\omega_0-\mathrm{i}\Gamma_r}+\frac{\langle f|\hat{p}_{\sigma}|r\rangle\langle r|\hat{p}_{\rho}|i\rangle}{\omega_{rf}+\omega_0+\mathrm{i}\Gamma_r}\right]$，
即分母\textbf{含阻尼项 $-\mathrm{i}\Gamma_r$}、系数为 $1/h$，\textbf{不含 $\nu_0^4$ 前因子}。
增强倍数原书印作 $10^4\sim10^6$ 倍（本分片已据此由 $10^2$—$10^6$ 改正）；首次实验观察到共振
拉曼效应者为 Shorygin（1953 年）。
\end{principle}

共振带来的最大好处是\textbf{选择性}：只有与那条电子跃迁耦合的模式被放大，因此可以专门去看骨
架杂原子（例如 TS-1 中与骨架钛相关的振动）、可以研究大分子聚集体的局部结构、也可以做共振拉
曼偏振测量从而得到正常拉曼效应中得不到的分子对称性信息。源文本列出的四条优点是：灵敏度高、
可检测低浓度和微量样品；通过共振拉曼谱带强度随激发线的关系可以给出分子振动和电子运动相互作
用的信息；偏振测量有时可给出分子对称性信息；利用标记分子基团的共振效应可研究大分子聚集体的
局部结构\src{S1 §6.4.2，书内 p.376 / PDF p.395}。

但选择性是双刃的，这一点必须记住：\textbf{既然不同模式的增强因子不同，共振拉曼谱的峰高比就
不再等于浓度比或模式数之比}。这与后面 SERS 的定量困境是同构问题——凡是有增强机制介入的拉曼
测量，绝对强度与峰高比都需要重新论证。

\subsubsection{紫外激发为什么同时避开荧光又提高灵敏度}
\label{sec:ch06b-uvmechanism}

这是本节的核心机制，值得拆成两问分别回答，因为两问的物理来源完全不同。

\textbf{第一问：为什么能避开荧光——能级间隙给出的波长分界。}物质分子吸收紫外光后，处于电子
基态的分子被激发跃迁到\textbf{一个高于第一电子激发态的虚态能级}，然后通过两种形式回到电子
基态：第一种从虚态能级直接跃迁回电子基态，即瑞利或拉曼散射过程；第二种从虚态能级通过\textbf{无
辐射跃迁}降到第一电子激发态的最低振转能级上，然后通过辐射跃迁回到电子基态，这是荧光或磷光
的辐射过程。\textbf{由于在虚态能级和第一电子激发态之间有一个间隙}，所以拉曼信号出现在比荧光
波长更短的区域，从而很有效地避开荧光干扰\src{S1 §6.4.1、§6.4.6，书内 p.374、p.382 /
PDF p.393、p.401}。

把这段话翻成可操作的判据，就是\textbf{能量账}：紫外激发把分子送到第一电子激发态之上的高虚
态；拉曼散射是"从哪来回哪去"，散射光子只比激发光子少一个振动量子，所以仍停在激发线附近的紫
外区；而荧光必须先把这个"高于第一电子激发态的间隙"以无辐射方式放掉，再发射，因此荧光光子
的能量固定地少了这一块，波长被推到长波侧。结果是：\textbf{紫外拉曼的谱线落在激发线附近几十
纳米以内，荧光带落在几十纳米之外}，两者在光谱上自然分开。相反，用可见光激发时，拉曼线与荧
光带挤在同一个波长区间里，只能靠扣基线把它们分开——这就是可见拉曼的根本困难。

\textbf{第二问：为什么能提高灵敏度——三项相加，不是一项。}

\begin{itemize}
\item \textbf{（i）$\nu_0^4$ 因子。}散射强度与激发光频率的四次方成正比。从 \qty{532}{nm} 换
  到 \qty{244}{nm}，频率因子提高 $(532/244)^4\approx 23$ 倍。这一项与共振无关，任何短波激发
  都享有。
\item \textbf{（ii）共振增强。}当紫外激光位于样品的吸收带时会产生共振效应，拉曼散射截面异常
  增大（上一小节），可以探测极低样品浓度的拉曼信号。这一项是最主要的贡献，量级可达
  $10^2$—$10^6$。
\item \textbf{（iii）把荧光从分母里删掉。}这一项最不显眼却常常最有效。信噪比的分母原本由荧光
  背底的涨落主导，换到紫外后分母变成散粒噪声主导；即使拉曼信号本身只增强一个数量级，可检性
  的提高也远不止一个数量级——因为原来是"在一个几万倍大的背底上找峰"。
\end{itemize}

\textbf{代价也要一起说清。}紫外拉曼并不是免费的：紫外激光与紫外光路昂贵（腔内倍频或
He—Cd 光源、紫外镀膜光栅与反射镜、紫外增强 CCD，见下一小节）；深紫外下样品吸收很强，探测深
度变浅，紫外拉曼因此变成\textbf{表面灵敏}技术——这既是它最大的方法论价值，也意味着\textbf{它
与可见拉曼测的可以不是同一层样品}（见后文 \ce{ZrO2}、\ce{TiO2} 相变的例子）；共振的选择性使
峰高比不能用于定量。至于"紫外激发是否同时降低激光加热"，不能一概而论——它取决于样品在紫外
区的吸收系数，必须用功率依赖实验实测。

图~\ref{fig:uv-fluor-model} 把上面三问的定量关系画在一张图上，并用它对照源文本给出的三条观测
（\qty{488}{nm} 只有强荧光背底、\qty{325}{nm} 有干扰但可得谱、\qty{244}{nm} 完全避开荧光）。

\begin{figure}[tbp]
\centering
\begin{tikzpicture}
% declare function 必须放在 \begin{tikzpicture} 之后、\begin{axis} 之前，用 \tikzset 声明。
% 荧光背底用逻辑斯蒂函数（有界指数的形式，指数绝对值 < 80，不会触发 Dimension too large）。
\tikzset{declare function={
  fluor(\x) = 1200/(1+exp(-(\x-318)/11));
  rnr(\x)   = (532/\x)^4;
  rres(\x)  = (532/\x)^4*(1+300/(1+((\x-244)/14)^2));
}}
\begin{axis}[
  width=\onecolwidth, height=64mm,
  xmin=220, xmax=700, ymin=1e-3, ymax=1e4,
  ymode=log,
  xlabel={激发波长 $\lambda$ / nm},
  ylabel={相对信号强度（对数标度）},
  tick label style={font={\fontsize{8}{9.6}\selectfont}},
  label style={font={\fontsize{8}{9.6}\selectfont}},
  legend style={at={(0.30,0.02)}, anchor=south west, draw=none,
     fill=white, fill opacity=0.9, text opacity=1,
     font={\fontsize{8}{9.6}\selectfont}, row sep=0.5pt, inner sep=1.5pt},
  every axis plot/.append style={thick, no markers, samples=250, domain=220:700},
]
\addplot[black] {fluor(x)};
\addlegendentry{荧光背底（模型）}
\addplot[black!60, dashed] {rnr(x)};
\addlegendentry{非共振拉曼（$\nu_0^4$）}
\addplot[black!30, dotted] {rres(x)};
\addlegendentry{共振增强拉曼（模型）}
\end{axis}
\end{tikzpicture}
\caption{模型计算：同一物种在不同激发波长下的拉曼相对强度与荧光背底对比\modelmark。读者应看
三件事：（i）荧光背底（实线）在约 \qty{320}{nm} 以上迅速升到 $1.2\times10^3$ 并在整个可见区
保持平坦，因此可见激发的困难不是"峰弱"而是"背底高"；（ii）非共振拉曼（虚线）只按 $\nu_0^4$
变化，从 \qty{532}{nm} 换到 \qty{244}{nm} 只提高约 23 倍；（iii）共振增强拉曼（点线）在
\qty{244}{nm} 出现共振峰，把荧光/拉曼比压到约 $1:4.7\times10^3$，而 \qty{325}{nm} 处约
$11:1$、\qty{488}{nm} 处约 $4\times10^2:1$——这正是"488 nm 只有很强的荧光背景、325 nm 有干
扰但得到了谱图、244 nm 完全避开荧光"三条观测的定量对应。曲线为模型计算：荧光背底取中心
\qty{318}{nm}、宽 \qty{11}{nm} 的逻辑斯蒂函数，共振项取中心 \qty{244}{nm}、宽 \qty{14}{nm}
的洛伦兹因子（峰值增强 $3\times10^2$ 倍），$\nu_0^4$ 因子以 \qty{532}{nm} 归一；参数是为说明
机制设定的，不代表任何一种具体样品的真实数值，三条比值也不应当作定量判据使用。
\figsrc{机制与三条观测据 S1 §6.4.1、§6.4.6.2，书内 p.374、p.386 / PDF p.393、p.405}}
\label{fig:uv-fluor-model}
\end{figure}

\begin{table}[htbp]
\centering
\caption{三条激发波长路线的机制与代价。读者应看第三列：三条路线不是"哪个更好"，而是各自把不
同的一项当作代价付出去；选路线等于选愿意付哪一项。\src{据 S1 §6.4.1，书内 p.374 / PDF p.393；
§6.4.3，书内 p.377—378 / PDF p.396—397；§6.4.6，书内 p.382—384 / PDF p.401—403 归纳。}}
\label{tab:raman-routes}
{\srcsize
\begin{tabularx}{\columnwidth}{@{}lXX@{}}
\toprule
路线 & 避开荧光的机制 & 付出的代价与边界 \\
\midrule
可见（常见 \qty{532}{nm}、\qty{488}{nm}） & 不避开，只能靠扣基线、时间分辨或样品前处理抑制（见本章 §6.2 的四条路线） & 荧光常比拉曼强几个数量级，有时根本得不到光谱；深颜色样品在可见区吸收强，既有荧光又有激光加热 \\
近红外（FT-Raman，常见 \qty{1064}{nm}） & 近红外光频率低，一般不会激发电子吸收带，分子的荧光不被激发 & 散射强度按 $\nu_0^4$ 下降（$532\to1064$ 降到 $1/16$）；红外探测器灵敏度比可见区光电倍增管差至少两个数量级；样品热辐射干扰大；稀土与过渡金属离子的近红外跃迁仍会带来荧光；原位温度上限低（$<\qty{150}{\degC}$） \\
紫外（\qty{325}{nm}、\qty{257}{nm}、\qty{244}{nm}） & 分子被激发到高于第一电子激发态的虚态，虚态与第一电子激发态之间有间隙，荧光必须先无辐射放掉这一块能量，因此拉曼信号出现在比荧光波长更短的区域 & 紫外光源与紫外光路昂贵；深紫外下样品吸收强，探测深度变浅而成为表面灵敏技术；共振使峰高比不再等比；深色样品在紫外区是否仍强吸收需实测 \\
\bottomrule
\end{tabularx}}
\end{table}

\emph{核对说明}：已对照原页面图像（书内 p.374 / PDF p.393）。原书 §6.4.1
的两条路线与本表「机制」列一致：（1）「将激发波长向近红外（大于 \qty{700}{nm}）
方向移动可以避开大部分荧光，因为近红外光频率低，一般不会激发电子吸收带（图 6-17），
很多研究者通过采用傅氏变换拉曼光谱得到了没有荧光干扰的拉曼光谱」；（2）「将激发波长向深紫外（小于
\qty{300}{nm}）
方向移动……处于电子基态的分子被激发跃迁到一个高于第一电子激发态的虚态能级……由于在虚态能级和第一电子激发态之间有一个间隙，
所以拉曼信号出现在比荧光波长更短的区域，从而很有效地避开荧光干扰」。可见区一行「不避开荧光」
与原书「荧光干扰和灵敏度较差是阻碍常规拉曼光谱得到广泛应用的最大问题」的判断一致。各「代价」
列（$\nu_{0}^{4}$ 使 532 换到 1064 nm 后散射强度降到约 1/16、
红外探测器性能比可见区差约 2 个数量级、原位温度上限低于 \qty{150}{\degC}、
紫外光源与光学元件昂贵且样品吸收强等）出自 §6.4.3（书内 p.377—378 / PDF p.396—397）
与 §6.4.6（书内 p.382—384 / PDF p.401—403）。本轮重核（PDF p.396—397、
p.402—403 = 书内 p.377—378、p.383—384 页面图像）：近红外一行的五条代价在书内
p.378 逐条可查——$\nu_0^4$ 使灵敏度大幅降低、红外区检测器比可见区光电倍增管灵敏度至少
差两个数量级、热辐射干扰较大、稀土与过渡金属离子的近红外电子跃迁仍无法避开荧光、样品
温度上限低于 \qty{150}{\degC}；紫外一行的四条代价（光源与光路昂贵、探测深度变浅而表面
灵敏、共振使峰高比不再等比、深色样品需实测）在源书只正面列优势，未见对应表述，如实记为
改编。另：「\qty{257}{nm}」源书印 257.2 nm（514.5 nm 倍频），本表按其常用写法取整。


\subsubsection{紫外拉曼谱仪的四个特殊之处}

源文本给出的紫外拉曼光谱仪结构（图~\ref{fig:uv-fluor-model} 之外的仪器部分）与常规可见拉曼
仪有四处不同，每一处都对应紫外波段的一个物理限制\src{S1 §6.4.6，书内 p.383—384 /
PDF p.402—403}。

\textbf{（1）光源：靠倍频与气体激光器得到单线紫外。}光源采用腔内倍频氩离子激光器与
He—Cd 激光器：腔内倍频的紫外激光器是把氩离子激光器的可见激光经腔内 BBO 晶体倍频，得到多条
单线紫外激光，常用的有 \qty{514.5}{nm} 倍频所得的 \qty{257}{nm} 和 \qty{488}{nm} 倍频所得
的 \qty{244}{nm} 激光线；He—Cd 激光器可以提供 \qty{325}{nm} 的激光。

\textbf{（2）外光路：椭球收集镜解决两个问题。}外光路系统采用 180° 背向散射方式进行信
号采集，收集镜采用椭圆收集镜——它同时解决了"常规透镜在紫外区难以使用"和"随激发波长变化需
要重新调节焦距"两个问题；背向散射信号经紫外镀膜的椭球镜反射收集、会聚到光谱仪的入射狭缝
进入分光系统。

\textbf{（3）光谱仪：三光栅色散相减去掉瑞利线。}光谱仪采用三光栅结构，为适应紫外波段，光栅
和反射镜均采用紫外区镀膜。前两块光栅处于镜像对称位置，进行光谱的\textbf{色散相减}，其作用是
利于中间狭缝滤掉瑞利线；去除瑞利线后的散射光经第三块光栅 G3 分光后照射到 CCD 探测器。这也
解释了为什么这台仪器要用三块光栅而不是单光栅：\textbf{低波数能力是靠色散相减换来的}，代价是
多一次色散带来的光通量下降（与本章仪器部件表中"多联光谱仪光通量下降"是同一条折中）。

\textbf{（4）探测器：为紫外增强的液氮冷却 CCD。}液氮冷却 CCD 采用背向探测和紫外增强，使其在
紫外区进行信号检测的量子效率达到 $50\%$ 以上；液氮冷却使暗电流极低（$1$—$3$ 电子/(像元$\cdot$小
时)），可以瞬时快速地采集信号，采集速率可达毫秒级。进一步发展的紫外光谱仪采用三联光谱仪结
构，每联在线配置三块光栅，可在紫外区到可见区（$200$—\qty{700}{nm}）的较宽范围内进行拉曼光
谱研究，分辨率在紫外区为 \qty{2.0}{cm^{-1}}、在可见区为 \qty{1.0}{cm^{-1}}。国内第一台连续
波紫外共振拉曼光谱仪由中国科学院大连化学物理研究所于 1997 年建成。

关于最后那个分辨率差异，有一条容易漏掉的换算关系：\textbf{波数分辨率与波长分辨率之间差一个
$1/\lambda^2$}，即 $\delta\tilde\nu=\delta\lambda/\lambda^2$。同一个波长色散量在更短的波长下
对应更大的波数间隔，因此紫外区的"$\mathrm{cm^{-1}}$ 分辨率"必须单独标注，不能直接与可见区
的数值比较。

\subsubsection{读图实例：0.1 wt\% \ce{MoO3}/$\gamma$-\ce{Al2O3} 在三条激发线下的谱图}
\label{sec:ch06b-mo-instance}

这一小节把"换激发线"的全部判断链走一遍。样品是负载量 \num{0.1} wt\% 的
\ce{MoO3}/$\gamma$-\ce{Al2O3}，源文本给出了它的紫外可见漫反射谱与三条激发线（\qty{488}{nm}、
\qty{325}{nm}、\qty{244}{nm}）下的拉曼谱\src{S1 §6.4.6.2，书内 p.386 / PDF p.405}。下面的
两栏\textbf{严格分开}：观察栏只写谱图上直接可读的事实，推断栏只写由这些事实出发的归属与结论；
任何一句话不允许在两栏之间来回借用。

\begin{obs}
（观察——谱图上直接可读，不含任何归属）
\begin{itemize}
\item 紫外可见漫反射谱上主要出现两个吸收峰，分别位于 \qty{220}{nm} 与 \qty{280}{nm}。
\item \qty{488}{nm} 激发：谱图中没有得到任何拉曼信号，只有很强的荧光背底。
\item \qty{325}{nm} 激发：存在荧光干扰，但得到了拉曼光谱图。
\item \qty{244}{nm} 激发：完全避开了荧光干扰，得到信噪比很好的拉曼光谱图。
\item \qty{244}{nm} 谱中可读出：\qty{325}{cm^{-1}} 与 \qty{910}{cm^{-1}} 两个峰，以及
  \qty{1802}{cm^{-1}} 与 \qty{2720}{cm^{-1}} 两个峰。
\item \qty{325}{nm} 谱中可读出：\qty{837}{cm^{-1}} 处的峰及其 \qty{1670}{cm^{-1}} 的倍频峰。
\end{itemize}
\end{obs}

\begin{infer}
（推断——归属与结论，每一条都指回上面某一条观察）
\begin{itemize}
\item \qty{220}{nm} 的吸收峰归属于\textbf{四配位}钼物种的吸收，\qty{280}{nm} 的吸收峰主要归
  属于\textbf{六配位}钼物种的吸收。
\item \qty{244}{nm} 的激发线位于 Mo—O 键的荷电跃迁区，共振效应使 Mo—O 键对称伸缩振动模的
  \textbf{基频、倍频和三倍频}的谱峰强度\textbf{同时}发生增强，因此才得到了传统拉曼光谱无法得
  到的倍频与三倍频峰。
\item \qty{325}{nm} 的激发线位于 Mo—O—Mo 桥键的紫外吸收带，因此 \qty{325}{nm} 谱中
  \qty{837}{cm^{-1}} 的峰与其 \qty{1670}{cm^{-1}} 的倍频峰对应桥键物种的\textbf{反对称}伸缩振动。
\item 因此选用 \qty{244}{nm} 与 \qty{325}{nm} 两条激发线，可以分别有选择地增强"端基 Mo—O"
  与"桥键 Mo—O—Mo"，从而有效区分表面钼物种的配位结构。
\item 结论一：在极低负载量（\num{0.1} wt\%）时，四配位和六配位物种\textbf{同时存在}。这与过
  去高负载量催化剂的研究结果不同——以前推测认为在低负载量时只有四配位钼存在。
\item 结论二：由于紫外共振拉曼增强效应，极低负载量氧化物催化剂的研究成为可能。
\end{itemize}
\end{infer}

\textbf{为什么必须把两栏分开。}"\qty{488}{nm} 下没有拉曼信号"是观察；"因为荧光把拉曼淹没
了"是推断，它依赖图~\ref{fig:uv-fluor-model} 的模型（模型给出 \qty{488}{nm} 处荧光/拉曼比约
$4\times10^2$）与源文本"只有很强的荧光背景"的记述，而不是谱图本身能读出的东西。同理，
"\qty{244}{nm} 谱中有 \qty{1802}{cm^{-1}} 的峰"是观察；"它是 \qty{910}{cm^{-1}} 基频的二倍
频"是推断，需要先做归一性检查。把这两类句子分开写，是判断"这条结论有多大把握"的最省力办法
——读别人的谱图时也一样。

\emph{核对说明}：已对照原页面图像（书内 p.386 / PDF p.405）。原书印：\qty{244}{nm} 激光激发的
拉曼光谱中出现了归属为四配位 Mo=O 键的 \qty{325}{cm^{-1}} 和 \qty{910}{cm^{-1}} 的弯曲及对称
伸缩振动峰，以及在 \qty{1802}{cm^{-1}} 和 \qty{2720}{cm^{-1}} 归属为 Mo=O 键对称伸缩振动的
二倍频和三倍频峰；\qty{325}{nm} 激发谱中可观察到位于 \qty{837}{cm^{-1}} 和 \qty{1670}{cm^{-1}}
的六配位 Mo—O—Mo \textbf{反对称}伸缩振动的基频和倍频峰（原稿误作"对称"伸缩，已改）；漫反射谱
两吸收峰为 \qty{220}{nm}（四配位氧化钼）与 \qty{280}{nm}（六配位氧化钼）。原书的这些波数
\textbf{彼此不自洽}：$2\times910=1820\neq1802$、$3\times910=2730\neq2720$、$2\times837=1674
\neq1670$，属\textbf{原书自身问题}；本分片按原书\textbf{照录观测波数}，不代为换算基频（原稿曾据
$2\times901=1802$ 反推基频为 \qty{901}{cm^{-1}}，经读图证伪，已改回原书的 \qty{910}{cm^{-1}}）。

\begin{dataflow}
从"换一条激发线"到"这条峰属于哪种配位"的完整链条是：\textbf{换线 $\to$ 看哪些峰改变 $\to$
判断改变来自共振还是来自荧光/吸收 $\to$ 只对增强的那几条峰做归属 $\to$ 用第二条落在不同吸收
带的激发线交叉验证 $\to$ 才落到"某种配位物种存在"}。每一步都有失效方式：
\begin{itemize}
\item \textbf{峰强改变也可能来自样品本身的变化}：激光加热、吸水或脱水都会改变谱峰（见
  \S\,\ref{sec:ch06b-insitu}）。换线实验应在同一批样品、同一预处理下尽快完成，并记录两次测
  量之间的间隔与样品存放条件。
\item \textbf{倍频峰必须先做归一性检查}：$\tilde\nu_2\approx2\tilde\nu_1$、
  $\tilde\nu_3\approx3\tilde\nu_1$。本实例正是靠这条检查发现基频波数的 OCR 疑点；不做检查就
  可能把两个不同物种的基频当成同一物种的一组倍频。
\item \textbf{"某条激发线有共振"这一步需要吸收谱作依据}：没有该样品的紫外可见漫反射谱，就不
  能说 \qty{244}{nm} 落在哪一条吸收带上，"共振"只能算猜测。归属与吸收带的对应关系（如
  \qty{220}{nm} 对四配位）同样来自文献与参考谱，不是谱图自动给出的。
\end{itemize}
\end{dataflow}

\begin{table}[htbp]
\centering
\caption{三条激发线各自给出的观测与可用判断（本实例）。读者应看最后一列：每条线能回答的问题
是有限的，把三条线合起来才能把"端基/桥键/四配位/六配位"这组问题问完。\src{据 S1 §6.4.6.2，
书内 p.385—387 / PDF p.404—406 归纳；表中标注「模型」的荧光/拉曼比值为本改编本模型计算，
非源书实验数据。}}
\label{tab:mo-lines}
{\srcsize
\begin{tabularx}{\columnwidth}{@{}lXX@{}}
\toprule
激发线 & 源文本给出的观测 & 由这条线能得到的判断（含代价） \\
\midrule
\qty{488}{nm}（可见） & 没有任何拉曼信号，只有很强的荧光背底 & 判断：该样品的可见激发路线不可用，必须换线。代价：$\nu_0^4$ 因子最小、荧光背底最高（模型约 $7\times10^2$ 倍），此时"测不到"本身不能推出"没有该物种" \\
\qty{325}{nm}（近紫外） & 有荧光干扰但得到谱图；约 \qty{837}{cm^{-1}} 的峰及其倍频（倍频波数 OCR 丢失） & 判断：该线位于 Mo—O—Mo 桥键吸收带，\qty{837}{cm^{-1}} 可作为桥键物种的判据。代价：荧光仍占优（模型约 $11$ 倍），且只有与桥键电子跃迁耦合的模式被增强 \\
\qty{244}{nm}（深紫外） & 完全避开荧光，信噪比好；\qty{910}{cm^{-1}} 处基频，以及 \qty{1802}{cm^{-1}} 与 \qty{2720}{cm^{-1}} 的倍频、三倍频 & 判断：该线位于端基 Mo—O 荷电跃迁带，共振使基频/倍频/三倍频同时增强，可测到 \num{0.1} wt\% 级负载量。代价：峰高比不等于物种比；深紫外探测深度浅，读到的主要是表层 \\
\bottomrule
\end{tabularx}}
\end{table}

\subsubsection{紫外拉曼在催化中的五个落点}

紫外拉曼之所以在催化研究中具有很大的优势和潜力，源文本归结为三点：灵敏度高、避开荧光干扰，
以及"由于很多化合物的电子吸收带在紫外区，还可以进行共振拉曼光谱的研究"
\src{S1 §6.4.6，书内 p.384 / PDF p.403}。下面五个落点是源文本给出实例最多的方向，也是读者
最可能直接照用的方向。

\textbf{（1）催化剂积炭失活。}催化剂表面的积炭物种主要是一些高度脱氢的含碳化合物，如烯烃、
稠环芳烃、石墨前驱体和石墨；这些物种的形成机理和表面状态很难研究。虽然拉曼光谱在理论上应该
是一种理想的表征表面积炭的技术，但由于碳氢化合物有很强的荧光干扰，很难用可见拉曼光谱进行表
征。积炭的谱峰主要出现在三个区域：$1360$—\qty{1400}{cm^{-1}}（C—H 变形振动）、$1580$—\qty{1640}{cm^{-1}}
（C=C 伸缩）与 $2900$—\qty{3100}{cm^{-1}}（C—H 伸缩振动）；通过对这些谱峰的位置和相对强度的分析，可以区分烯烃、聚烯烃、芳烃、聚芳烃、类
石墨等不同形态的积炭。采用紫外激发线，不但使拉曼散射截面增加，而且有效避开荧光干扰，得到信
噪比很好的紫外拉曼光谱\src{S1 §6.4.6.1，书内 p.384 / PDF p.403}。

三种分子筛在甲醇转化反应中的积炭物种不同，这是"同类技术给出差异"的一个漂亮例子：
\begin{itemize}
\item \textbf{SAPO-34}：紫外拉曼谱中出现 1414、1616、1628、2822 与 \qty{2974}{cm^{-1}} 五个谱峰。
  其中 \qty{1414}{cm^{-1}} 谱峰为 \ce{CH3} 的变形振动，1616 与 \qty{1628}{cm^{-1}} 为烯烃的 C=C
  伸缩振动，2822 与 \qty{2974}{cm^{-1}} 为 C—H 键的对称与反对称伸缩振动——说明 SAPO-34 中形成
  的积炭物种主要是由\textbf{富含氢的烯烃和聚烯烃分子}组成。
\item \textbf{ZSM-5}：紫外拉曼谱中 \qty{1425}{cm^{-1}} 与 \qty{1615}{cm^{-1}} 的谱峰分别归属为
  芳烃和取代芳烃的 C—H 变形振动与 C=C 伸缩振动；在 ZSM-5 与 USY 的紫外拉曼谱图中\textbf{未检
  测到 C—H 振动区间的谱峰}。
\item \textbf{USY}：\qty{1604}{cm^{-1}} 谱峰为聚芳烃的 C=C 伸缩振动。
\end{itemize}
源文本的结论是：SAPO-34、ZSM-5 与 USY 三种分子筛在甲醇转化反应中形成了不同的积炭物种，这
是由它们不同的\textbf{酸性和孔结构}所决定的\src{S1 §6.4.6.1，书内 p.384—385 /
PDF p.403—404}。

\emph{核对说明}：已对照原页面图像（书内 p.384—385 / PDF p.403—404）。SAPO-34 的五个谱峰原书
印作 1414、1616、1628、2822 与 \qty{2974}{cm^{-1}}：\qty{1414}{cm^{-1}} 为 \ce{CH3} 的变形振
动，1616 与 \qty{1628}{cm^{-1}} 为烯烃的 C=C 伸缩振动，2822 与 \qty{2974}{cm^{-1}} 为 C—H 键
的对称与反对称伸缩振动；图 6-26 的 SAPO-34 曲线亦标注 2822 与 2974。原书归属句把后者印作
\qty{2947}{cm^{-1}}，与同页峰列表及图 6-26 标注不一致，属\textbf{原书自身问题}，本分片按峰列表
与图注取 2974。ZSM-5 与 USY 的归属亦据同页改正：ZSM-5 的 1425 与 \qty{1615}{cm^{-1}} 分别归属
为\textbf{芳烃和取代芳烃}的 C—H 变形振动与 C=C 伸缩振动（原稿误作"烯烃……C—C"），USY 的
\qty{1604}{cm^{-1}} 为聚芳烃的 C=C 伸缩振动（原稿漏数值）；"ZSM-5 与 USY 未检测到 C—H 振动区
间的谱峰"一句原书印于 SAPO-34 段末，本分片照录。积炭谱峰的三区间（$1360$—\qty{1400}{cm^{-1}}
C—H 变形、$1580$—\qty{1640}{cm^{-1}} C=C 伸缩、$2900$—\qty{3100}{cm^{-1}} C—H 伸缩）亦据同页
补正（原稿漏后两区间的数值）。

\textbf{（2）低负载量负载型过渡金属氧化物催化剂。}由于过渡金属氧化物在紫外区有荷电跃迁吸收，
选择合适的紫外激光进行共振激发，可以很有效地得到低负载量催化剂的紫外拉曼光谱。李灿等把紫外
拉曼用于低负载量 \ce{MoO3}/$\gamma$-\ce{Al2O3} 的研究，结果表明由于避开了荧光干扰、激发波长
变短和共振效应，信号灵敏度大幅度提高，甚至可以得到负载量低于 \num{0.1} wt\% 的
\ce{MoO3}/$\gamma$-\ce{Al2O3} 催化剂的谱图（即上一小节的读图实例）\src{S1 §6.4.6.2，书内
p.385—386 / PDF p.404—405}。

更重要的是这条路线能跟随\textbf{制备过程}：源文本报道用紫外拉曼研究负载型 \ce{MoO3} 催化剂
的制备与焙烧过程，结果表明在湿润状态时表面物种的配位结构不仅与浸渍液的 pH 有关，而且与载体
的性质有关；焙烧后几乎所有催化剂的表面钼物种都发生了变化，导致\textbf{四配位物种减少、六配
位物种增加}。作者用紫外拉曼灵敏地检测到负载钼物种随载体表面 pH 改变而发生的配位状态变化，
发现\textbf{表面 pH 是决定表面钼物种配位结构的重要因素}。由于绝大部分过渡金属氧化物在紫外区
有吸收，这项研究可以推广到其他氧化物\src{S1 §6.4.6.2，书内 p.387 / PDF p.406}。最后一句的方
法论含义值得单独记住：\textbf{紫外拉曼给出的不是"最终催化剂长什么样"，而是"制备条件如何一步
步写进表面配位结构"}。

\textbf{（3）原位跟踪分子筛的水热合成。}分子筛大多在高温高压的水热体系中合成，把中间产物取
出反应体系可能导致其结构性质发生明显变化；而原位表征可以在合成过程中测定中间物种，连续全面
地跟踪反应过程，得到最真实可靠的信息\src{S1 §6.4.6.3，书内 p.388 / PDF p.407}。拉曼在这里
有一个别人没有的条件：\textbf{水的拉曼散射截面非常小}，所以可以同时研究液相与固相信号；而合
成体系中通常存在大量富含模板剂的物种，它们产生极强的荧光，极大干扰常规拉曼测量——紫外拉曼
正好能避开这些中间物种的荧光干扰\src{S1 §6.4.6.3，书内 p.388 / PDF p.407}。原位池的做法见
\S\,\ref{sec:ch06b-insitu}。

X 型分子筛（\qty{373}{K}，水热条件）的两组观测构成了完整的机理证据链：
\begin{itemize}
\item \textbf{液相的观测}：最初的谱图有一个中心位于 \qty{500}{cm^{-1}} 的宽峰，以及位于 774
  与 \qty{920}{cm^{-1}} 的比较明显的谱峰。\qty{500}{cm^{-1}} 附近的谱峰归属于无定形硅铝凝
  胶前驱体；\qty{1060}{cm^{-1}} 的谱峰则是合成过程中碱性凝胶吸收空气中的 \ce{CO2} 产生的物
  种；774 与 \qty{920}{cm^{-1}} 被归属为液相中的\textbf{单态硅物种}。\qty{774}{cm^{-1}} 谱
  峰的强度随温度升高而逐渐加强，晶化到 \qty{120}{min} 时达到最大值，随后随反应时间延长逐渐
  下降——说明合成过程中存在大量单态硅物种，并且\textbf{单态硅物种参与了整个晶化过程}。
\item \textbf{由此描述的过程}：最初阶段无定形硅凝胶与液相中的物种处于平衡状态；随着反应温度
  提高，大量单态硅物种从硅凝胶前驱体当中解聚出来；随着反应进行，晶体的生长需要更多的单态硅
  物种；到晶化后期几乎所有无定形前驱体都被消耗并转变为晶形结构，不再解聚出单态硅物种。
\item \textbf{固相的观测}：初始凝胶处于无定形状态，给出位于 500、575、920 与 \qty{1060}{cm^{-1}}
  的谱峰。\qty{575}{cm^{-1}} 的谱峰在反应刚开始就存在，并且随着反应进行强度逐渐下降；当谱图
  中出现 298 与 \qty{380}{cm^{-1}} 的谱峰时，\qty{575}{cm^{-1}} 的峰基本消失——说明它代表合
  成中的一种\textbf{关键物种}。$550$—\qty{575}{cm^{-1}} 的谱峰被归属为 Al—O—Si 的伸缩振动；
  由于 \qty{575}{cm^{-1}} 与四元环呼吸振动的位置相差很远，它代表的 Al—O—Si 可能是以\textbf{支
  链}的形式存在，即来自位于四元环支链的 Al—O—Si 弯曲振动。
\item \textbf{晶化完成的判据}：随反应进行，固相 \qty{500}{cm^{-1}} 的谱峰移至
  \qty{514}{cm^{-1}}，后者归属于晶化完全的 X 型分子筛；同时对应于 X 型分子筛的特征峰
  \qty{298}{cm^{-1}} 与 \qty{380}{cm^{-1}} 强度增强，它们是由分子筛结构中\textbf{四元环和双
  六元环}的呼吸振动产生的。此外，固相中 \qty{774}{cm^{-1}} 与 \qty{920}{cm^{-1}} 谱峰的存
  在说明一部分单态硅物种吸附在固相表面或者是陷于固相中间。
\end{itemize}
综合原位拉曼、理论计算与 XRD、核磁共振谱，源文本给出 X 型分子筛合成机理的图景：无定形硅铝凝
胶不断解聚、溶解形成可溶性单态硅物种（\qty{774}{cm^{-1}}）；前驱体中存在四元环结构单元，其
中一些以支链形式存在（\qty{575}{cm^{-1}}）；四元环结构单元之间相互连接形成部分结晶的分子筛
结构（\qty{500}{cm^{-1}}）；晶化完全后该峰移至 \qty{514}{cm^{-1}}，并出现 \qty{298}{cm^{-1}}
与 \qty{380}{cm^{-1}} 的呼吸振动峰\src{S1 §6.4.6.3，书内 p.388—392 / PDF p.407—411}。

\textbf{（4）杂原子分子筛：共振选择性给出"杂原子是否进入骨架"的直接证据。}大部分分子筛催化
剂有较强的荧光干扰问题，原因之一是制备分子筛时残留的有机模板剂和合成原料中的杂质引起荧光；
因此在进行可见拉曼光谱表征前往往需要长时间的焙烧来消除荧光干扰，即使如此，很多分子筛的可见
拉曼光谱表征也是很困难的\src{S1 §6.4.6.4，书内 p.392 / PDF p.411}。三个实例：
\begin{itemize}
\item \textbf{TS-1}：以 \qty{244}{nm} 紫外激光作为激发光进行表征，通过紫外激光有选择性地激
  发了 TS-1 分子筛中 Ti—O—Si 之间的荷电跃迁，使与\textbf{骨架钛}直接相关的拉曼峰共振增强，
  得到钛进入骨架的最直接证据。TS-1 的紫外可见漫反射光谱在 \qty{220}{nm} 处有吸收峰，该峰归
  属于 Ti—O—Si 键骨架之间的 $\mathrm{p}\pi$—$\mathrm{d}\pi$ 荷电跃迁，而 Silicalite-1 在紫
  外区没有该吸收。与可见拉曼相比，紫外拉曼出现了 490、530 与 \qty{1125}{cm^{-1}} 三个新峰，
  这是骨架钛和骨架晶格氧原子之间的电荷跃迁受到紫外激光激发后产生的共振拉曼峰；由于共振效应，
  与骨架钛物种相关的拉曼谱峰增加了几个数量级。
\item \textbf{V-MCM-41}：紫外拉曼谱中 \qty{490}{cm^{-1}} 谱峰变宽，并在 930 与
  \qty{1070}{cm^{-1}} 出现两个新峰。\qty{930}{cm^{-1}} 的谱峰为\textbf{骨架外聚合}氧化钒的 V=O 对称伸缩振动峰，
  \qty{1070}{cm^{-1}} 的谱峰则为骨架四配位氧化钒的 V=O 对称伸缩振动峰。作者认为在可见拉曼光谱中未发现这两
  个峰，是由于 \qty{244}{nm} 波长的激发线\textbf{同时}激发了骨架钒和非骨架钒物种的荷电跃迁，
  共振效应使这两个峰的强度增强，从而同时得到了骨架钒与非骨架钒物种的谱峰。
\item \textbf{Fe-ZSM-5}（Si/Fe $=76$）：紫外可见漫反射谱最特征的吸收带位于 211 与
  \qty{245}{nm}，被归属为分子筛中骨架氧到骨架铁原子的荷电跃迁。三条激发线对比：\qty{244}{nm}
  位于该荷电跃迁吸收带、\qty{325}{nm} 处于它的吸收带边、\qty{532}{nm} 远离吸收带。\qty{244}{nm}
  谱中除 290、380 与 $800$ 附近的分子筛骨架特征峰以外，还激发出 516、1005、1115 与
  \qty{1165}{cm^{-1}} 四个新峰（这些峰在纯 ZSM-5 中是没有的）；\qty{325}{nm} 谱中这些峰仍
  在但强度弱了很多；\qty{532}{nm} 谱中只有 \qty{1005}{cm^{-1}} 的谱峰。作者把峰分为两类：
  \qty{1005}{cm^{-1}} 由铁的引入引起，但它是一个\textbf{非共振}峰；据对钛硅分子筛的归属，516 与
  \qty{1165}{cm^{-1}} 可归属为对称与反对称的 \ce{FeOSi} 四面体振动模式（原书未单独归属 1115）。此外，利用紫外拉曼还检测到了
  ZSM-5 分子筛中\textbf{痕量铁}的存在。
\end{itemize}
V-MCM-41 的例子里藏着一个重要提醒：\textbf{共振选择性的前提是"只有目标跃迁落在激发窗口内"}。
一旦两条跃迁（骨架钒与非骨架钒）落在同一波长窗口，选择性就消失——此时紫外拉曼把两类物种一起
增强，反而无法区分它们。

\emph{核对说明}：已对照原页面图像（书内 p.393—394 / PDF p.412—413）。图 6-35 题注为「Fe-ZSM-5
（Si/Fe $=76$）的（a）紫外可见漫反射谱图；（b）\qty{244}{nm}；（c）\qty{325}{nm} 和（d）
\qty{532}{nm} 激发的拉曼光谱图」。正文印：紫外可见吸收光谱图中最特征的吸收带位于
\textbf{\qty{211}{nm} 和 \qty{245}{nm}}，归为分子筛中骨架氧到骨架位铁原子的荷电跃迁；\qty{244}{nm}
谱中「除了位于 \qty{290}{cm^{-1}}、\qty{380}{cm^{-1}} 和 \qty{800}{cm^{-1}} 等分子筛骨架的特征峰
以外」还激发出 516、1005、1115 与 \qty{1165}{cm^{-1}} 四个新峰（原稿「$800$ 附近」已按原书改为
精确值）；\qty{325}{nm} 谱中 \qty{1005}{cm^{-1}} 仍见而 516/1115/1165 强度弱了很多；\qty{532}{nm}
谱中只有 \qty{1005}{cm^{-1}}。峰的两类划分亦据原书：\qty{1005}{cm^{-1}} 由铁的引入引起但是
\textbf{非共振}峰，「根据对钛硅分子筛的归属，\qty{516}{cm^{-1}} 和 \qty{1165}{cm^{-1}} 可以归属为
对称和反对称的 \ce{FeOSi} 四面体的振动模式」——\qty{1115}{cm^{-1}} 原书\textbf{未单独归属}，本分片
不再把它列入四面体振动模式。

\textbf{（5）氧化物表面相变：紫外拉曼为什么比 XRD 更"表面"。}这是本节方法论价值最高的一段，
它把"信息深度"从一个抽象概念变成了可验证的读数差异。

\textbf{判据的物理基础}：一般来讲，拉曼散射信号是同时来自样品体相和表面的；但当样品对激发
和散射光有很强吸收时，来自体相的信号会很大程度地衰减，拉曼谱图反映的主要是\textbf{样品表面
区}的信息。大多数过渡金属氧化物对紫外光有很强的吸收，因而紫外拉曼光谱对这些样品的表面信息
比体相信息更为灵敏；而大多数样品在可见区都没有吸收，因此可见拉曼光谱给出的是\textbf{体相和
表面二者混合}的信息\src{S1 §6.4.6.5，书内 p.395 / PDF p.414}。

\begin{itemize}
\item \textbf{\ce{ZrO2}：表面先相变。}紫外拉曼谱显示：\qty{400}{\degC} 焙烧后可以观察到四方
  和单斜混合相的拉曼谱峰；\qty{500}{\degC} 焙烧后四方相谱峰完全消失，仅有单斜相的谱峰被观察
  到；升温到 \qty{700}{\degC}，单斜相谱峰进一步增强。但 XRD 图谱显示 \qty{400}{\degC} 焙烧
  后氧化锆主要以\textbf{四方相}结构存在，甚至焙烧样品在 \qty{700}{\degC} 后仍然能观察到四方
  相的存在；可见拉曼与 XRD 的结果非常相似，而与紫外拉曼的结果却明显不同。作者的结论是：氧化
  锆的相变是一个从表面到体相的过程，即\textbf{单斜相首先在四方相晶粒的表面形成}；尽管单斜相
  是热力学更稳定的相，但由于动力学的原因，从四方相到单斜相的相变需要一个缓慢的从表面到体相
  的过程。
\item \textbf{掺杂 \ce{ZrO2}：体相稳定不等于表面稳定。}从紫外拉曼、XRD 和可见拉曼的结果看，
  尽管由于稳定剂的加入，四方相能稳定存在于样品体相，但\textbf{四方相在表面仍不能被稳定}，
  极易转变为单斜相；特别是当稳定剂含量低时，样品表面完全处于单斜相结构；稳定剂在样品表面含
  量足够高时，样品可完全稳定在四方相。
\item \textbf{\ce{TiO2}：表面后相变，方向与 \ce{ZrO2} 相反。}XRD 与可见拉曼给出的是体相信
  息：焙烧温度达到 \qty{550}{\degC} 时，除锐钛矿的特征衍射峰之外还在 $27.6$°、$36.1$°、
  $41.2$° 和 $54.3$° 出现金红石的特征衍射峰，表明此时发生了从锐钛矿到金红石的相变；可见拉曼
  在 \qty{550}{\degC} 观察到很弱的金红石特征峰（\qty{445}{cm^{-1}}），\qty{580}{\degC} 之后
  又观察到金红石的另外两个特征峰 \qty{235}{cm^{-1}} 和 \qty{612}{cm^{-1}}；温度达到
  \qty{750}{\degC} 时锐钛矿的特征峰完全消失，相变完成。紫外拉曼给出的表面信息不同：样品在
  \qty{500}{\degC} 焙烧后可以明显观察到锐钛矿的特征峰（143、395、\qty{515}{cm^{-1}}），直到
  温度升高到 \qty{700}{\degC} 才能在紫外拉曼光谱上观察到金红石的特征峰（\qty{235}{cm^{-1}}）。
  因此\textbf{表面上锐钛矿的相变要滞后一些}，作者认为氧化钛的相变是从\textbf{团聚粒子的体相
  开始、逐渐发展到表面区域}的。
\item \textbf{把表面晶相接到催化活性上。}采用光催化分解水作为指标反应，考察不同温度焙烧的氧
  化钛催化剂的产氢活性：只要表面处于锐钛矿相（焙烧温度 $<\qty{680}{\degC}$），即使体相的锐钛
  矿—金红石含量有很大差别，样品的产氢活性基本没有明显变化；\qty{700}{\degC} 和
  \qty{750}{\degC} 焙烧之后光催化活性有所提高——这两个样品的体相中基本上均处于单一金红石晶
  相，但\textbf{表面处于锐钛矿和金红石的混合晶相}；焙烧温度升高到 \qty{800}{\degC} 时表面完全
  处于金红石晶相，产氢活性急剧降低。作者认为表面锐钛矿—金红石的\textbf{异相结}有利于电子—
  空穴分离，从而提高光催化活性。为验证这一结论，他们用浸渍并焙烧的方法制备了金红石负载的锐
  钛矿样品 \ce{TiO2(A)/TiO2(R)}-$n$（$n$ 为浸渍次数）：浸渍次数从 1 增加到 4 次时，用紫外拉
  曼的估算结果表明表面锐钛矿含量从 $10\%$ 增加到 $36\%$；\ce{TiO2(A)/TiO2(R)}-3 样品的光催
  化活性基本上是纯金红石载体的 4 倍，而未经过焙烧的样品活性非常低（甚至比金红石载体还低，可
  能因为金红石表面上覆盖了一层无定形 \ce{TiO2}）。上述结果说明：\textbf{相比于体相晶相，光催
  化活性更多地取决于表面区的晶相}。
\item \textbf{钼酸铁：表面富集。}\ce{Fe2(MoO4)3} 也是一类在紫外区有强吸收的材料。可见拉曼谱
  随焙烧温度升高，\ce{Fe2(MoO4)3} 与 \ce{MoO3} 的特征峰相对强度变化不大；紫外拉曼谱则表明随
  焙烧温度升高 \ce{MoO3} 的特征峰强度不断增加，\qty{700}{\degC} 焙烧后主要显示为 \ce{MoO3}
  的特征峰——说明随焙烧温度升高 \ce{MoO3} 逐步\textbf{在表面富集}。扫描电镜的结果支持这一判
  断：\qty{500}{\degC} 焙烧的样品主要为球形颗粒的聚集体，\qty{700}{\degC} 焙烧后除颗粒长大
  外还可明显观察到棒状物种出现；文献一般认为这种形貌为 \ce{MoO3} 的特征形貌，而且 EDX 结果
  也表明棒状区域为富钼区域。
\end{itemize}
源文本对这一段的总评是：紫外拉曼不仅对研究如氧化锆、氧化钛及钼酸铁等系列在紫外区有强吸收的
材料的表面相变及表面相组成具有优势，而且对研究半导体材料的光催化性能具有很大的优势
\src{S1 §6.4.6.5，书内 p.396—402 / PDF p.415—421}。

\begin{table}[htbp]
\centering
\caption{紫外拉曼在催化中的五个落点：关键谱峰与"为什么必须用紫外"。读者应看最后一列——每个
落点都需要紫外，但需要的\textbf{理由并不相同}：前四个靠"避开荧光 $+$ 共振选择性"，第五个靠
"强吸收导致的表面灵敏"。\src{据 S1 §6.4.6.1—§6.4.6.5，书内 p.384—402 / PDF p.403—421 归纳。}}
\label{tab:uv-apps}
{\srcsize
\begin{tabularx}{\columnwidth}{@{}lXX@{}}
\toprule
应用对象 & 紫外拉曼给出的关键谱峰 & 为什么必须用紫外 \\
\midrule
分子筛积炭（SAPO-34、ZSM-5、USY） & 积炭总体在 $1360$—\qty{1400}{cm^{-1}}；SAPO-34：1414、1616、1628、2822、2974；ZSM-5：1425、1615；USY：1604（聚芳烃 C=C 伸缩） & 碳氢化合物荧光极强，可见拉曼很难表征；紫外激发同时提高散射截面并避开荧光 \\
低负载量 \ce{MoO3}/$\gamma$-\ce{Al2O3}（\num{0.1} wt\%） & \qty{244}{nm}：910、1802、2720；\qty{325}{nm}：837 及倍频 & 激发线落在 Mo—O 与 Mo—O—Mo 的荷电跃迁带内，共振使信号增强数个数量级；可见激发只剩荧光 \\
X 型分子筛水热合成（原位） & 液相：500、774、920、1060；固相：500（晶化后移至 514）、575、298、380 & 合成体系含大量有机模板剂，荧光极强；水的拉曼截面很小，故可同时测液相与固相 \\
杂原子分子筛（TS-1、V-MCM-41、Fe-ZSM-5） & TS-1：490、530、1125；V-MCM-41：930、1070；Fe-ZSM-5：516、1005、1115、1165 & 需要\textbf{选择性}共振激发杂原子—氧的荷电跃迁，从而只增强与骨架杂原子相关的模式 \\
氧化物表面相变与表面组成（\ce{ZrO2}、\ce{TiO2}、\ce{Fe2(MoO4)3}） & \ce{TiO2} 锐钛矿 143、395、515 / 金红石 235、445、612；\ce{ZrO2} 四方与单斜各自的谱峰 & 紫外强吸收把体相信号衰减掉，信号主要来自表面区，从而把\textbf{表面相}与\textbf{体相}分开 \\
\bottomrule
\end{tabularx}}
\end{table}

\emph{核对说明}：本表第一行「分子筛积炭」已对照原页面图像（书内 p.384—385 / PDF p.403—404），
原书 SAPO-34 的谱峰为 1414、1616、1628、2822 与 \qty{2974}{cm^{-1}}，ZSM-5 为 1425 与
\qty{1615}{cm^{-1}}，USY 为 \qty{1604}{cm^{-1}}（聚芳烃的 C=C 伸缩振动）；本表原印 2947 与
「C—C 伸缩」均与原文不符，已改正。本轮重核其余四行（PDF p.405、p.409、p.411、p.415 = 书内 p.386、p.390、p.392、p.396 页面图像）：
低负载 \ce{MoO3}/$\gamma$-\ce{Al2O3} 一行与书内 p.386 的 244 nm（325、910、1802、\qty{2720}{cm^{-1}}）
与 325 nm（\qty{837}{cm^{-1}} 及倍频）两组波数一致；X 型分子筛一行与书内 p.390 的液相
500、774、920、\qty{1060}{cm^{-1}} 与固相 500（晶化后移至 514）、575（肩峰）、298、
\qty{380}{cm^{-1}} 逐项一致；杂原子分子筛一行 TS-1 的 490、530、\qty{1125}{cm^{-1}} 与
p.392 逐字一致（V-MCM-41 与 Fe-ZSM-5 两组在同节后页，未逐字对读，本表如实标为差异）；
氧化物表面相变一行的「自吸收使信号只来自表面」机制与 p.396 逐字一致，锐钛矿/金红石的六个
峰位未在所读页面出现，如实标为差异。另记源书自身的一处不一致：X 型分子筛的固相峰在 p.390
正文作 298、在 p.392 图 6-33(b) 标注作 290，本表按正文取 298。

\subsection{表面增强拉曼与共焦拉曼}

\subsubsection{SERS 的两种增强机制}

自 1974 年 Fleischmann 等首次从吸附物种的拉曼光谱中发现表面增强拉曼散射
（Surface-Enhanced Raman Scattering，SERS）以来，这方面工作发展非常迅速，理论和实验研究均
有大量报道，目前已成为拉曼光谱中一个非常活跃的领域\src{S1 §6.4.4，书内 p.378 / PDF p.397}。
SERS 的价值在于它把"极弱"的拉曼散射放大到可以测单层吸附物种的程度——这正是催化表面研究最
想要的能力。

\begin{principle}
\textbf{两种机制，对应诱导偶极公式的两个因子。}关于 SERS 的机理基本上有两种观点，把
$P=\alpha E$ 写开就一目了然：

\begin{equation}
P \;=\; \alpha\,E:\qquad
\text{电磁机制使 } E \text{ 增大},\qquad
\text{化学机制使 } \alpha \text{ 增大}
\label{eq:sers-two}
\end{equation}

\textbf{电磁机制}（物理增强为主）：认为 SERS 起源于激发光和拉曼散射光被粗糙的金属表面局部
电场增强，即 $P=\alpha E$ 中的 $E$ 增大。\textbf{化学机制}（化学吸附为主）：认为增强来自吸
附分子与金属之间的电荷转移而引起的分子极化率增加，即 $P=\alpha E$ 中的 $\alpha$ 增大。

\textbf{两种机制对定量的含义完全不同}：电磁增强几乎不挑分子种类（挑的是分子离金属表面的距
离与取向，以及它是否落在局域电场增强区里），因此它对同一分子内不同振动模式的影响相对温和，
但随位置剧烈变化；化学增强只出现在能形成特定吸附键与电荷转移共振条件的分子—金属组合上，并
会\textbf{改变相对峰高}。真实体系的净增强同时含有两个贡献，无法用一组标定把它们分开
\src{S1 §6.4.4，书内 p.378 / PDF p.397}。
\end{principle}

能产生 SERS 效应的金属很多，如 Ag、Au、Cu、Fe、Co、Ni、Al 等，常用的是 Ag、Au 和 Cu；获得
活性金属表面的方法有溶胶化学还原法、真空蒸馏法与化学沉积法\src{S1 §6.4.4，书内 p.378 /
PDF p.397}。源文本也直接给出了这条技术的边界：利用 SERS 研究\textbf{所有的金属和金属氧化物}
表面增强是不可能的，研究最多的是 Ag、Au 和 Cu；在这三种常用基质上可以进行有机分子、
\ce{SO2}、\ce{NO2}、\ce{CO}、\ce{NH3} 等物质的吸附和反应研究\src{S1 §6.4.4，书内 p.378 /
PDF p.397}。

三个有代表性的实验事实：
\begin{itemize}
\item \textbf{共吸附会改变基底。}有作者用 SERS 研究多孔 Ag 薄膜，观察到了 \ce{N2} 和 \ce{CO}
  增强的拉曼谱带，但 \ce{O2} 的拉曼谱带并没有被明显增强；若在吸附吡啶和乙烯以前，先在 Ag 表
  面吸附少量 \ce{O2}，吡啶和乙烯吸附物种的增强谱带会很大衰减。作者认为这可能是 \ce{O2} 的存
  在减少了其他分子在表面的吸附浓度、引起了其他分子的解离、改变了表面形态。注意作者只给出
  "可能"的解释——这本身就是一条读法提示：\textbf{预吸附改变了基底本身}，此后任何强度比较都不
  再可比。
  \item \textbf{增强因子的量级。}Feilchenfeld 和 Siiman 研究了 \ce{CrO4^{2-}} 在 Ag 溶胶
  中吸附的 SERS，其中 \textbf{Cr—O 键}的伸缩振动模式最大增强接近 $\num{1.1e5}$，此时激发
  光波长为 \qty{580.0}{nm}
  \src{S1 §6.4.4，书内 p.378—379 / PDF p.397—398}。
\item \textbf{电位依赖的 C—O 伸缩频率。}Tadayyoni 和 Weaver 报道了 \ce{HClO4} 溶液中 CO 的
  SERS：C—O 伸缩振动频率随电极上电压的改变在 $2080$—\qty{2110}{cm^{-1}} 之间变化；\textbf{不加
  电压}时，C—O 振动频率位于 \qty{2015}{cm^{-1}}——这与 CO 在催化剂表面的吸附频率完全一致。
  SERS 真正的强项：它测的是\textbf{吸附态的键与取向}，而且可以在电化学条件下原位测。
\end{itemize}

\emph{核对说明}：已对照原页面图像（书内 p.378—379 / PDF p.397—398）。原书印「Cr—O 键伸缩振动
模式最大增强接近 $1.1\times10^{5}$，此时激发光波长为 \qty{580.0}{nm}」——原稿因 OCR 损坏
（作「ｃ卜。键伸缩振动」）只写作"某一键的伸缩振动模式"，现按原书补出\textbf{Cr—O 键}。该例为
Feilchenfeld 和 Siiman$^{[165]}$ 研究 \ce{CrO4^{2-}} 在 Ag 溶胶中吸附的 SERS，图 6-21 为吸附的
\ce{CrO4^{2-}} 随激发光频率变化的表面增强拉曼光谱（激发波长 457.9、488.0、514.5、593.1、607.2、
676.4 nm，图上另标 \qty{796}{cm^{-1}} 的 \ce{CrO4^{2-}} on Ag 与 \qty{1020}{cm^{-1}} 的
\ce{CH3OH}）。C—O 一例原书印「Tadayyani 和 Weaver$^{[166]}$ 报道了在 \ce{HClO4} 溶液存在下 Au
电极上 CO 吸附的 SERS。C—O 伸缩振动频率随电极上电压的改变在 $2080\sim2110\,\mathrm{cm^{-1}}$
之间变化。\textbf{没有电压时}，C—O 振动频率在 \qty{2015}{cm^{-1}}，这与 CO 在催化剂表面的吸附
频率是完全一致的」——原稿误作"某一电压下"，已改为"不加电压时"；第二个姓原书自身印作
\textbf{Tadayyani}，本分片按通行文献拼写取 Tadayyoni。

\subsubsection{SERS 为什么不给浓度}
\label{sec:ch06b-sers-quant}

\begin{pitfall}
\textbf{把 SERS 的峰强度直接当作浓度}，是这项技术最典型、也最容易被"极高灵敏度"这句话诱导
出来的误判。三条理由互相独立，任何一条都足以推翻定量结论。

\textbf{（1）热点分布决定强度，而热点是纳米尺度的。}电磁增强来自粗糙金属表面（颗粒间隙、
尖端、聚集体缝隙）的局域电场增强，有效区域只占表面的很小一部分。分子落在热点里与落在热点外，
增强因子可以相差若干数量级，因此谱峰强度反映的是"\textbf{分子恰好位于热点中的概率} $\times$
表面浓度"，而前者随基底形貌、颗粒尺寸乃至分子自身大小变化。把强度当作浓度的直接读数，等于
假设热点分布对所有样品相同。

\textbf{（2）吸附取向改变相对峰高。}式~\eqref{eq:sers-two} 中的 $\alpha$ 是张量，SERS 强度
取决于分子某一振动模式的极化率张量分量与表面局域电场方向的匹配。同一个分子以不同取向吸附时，
不同振动模式的增强因子不同，所以\textbf{不仅绝对强度不可比，峰高比也不可比}；比较两个样品的
同一峰高比，等于默认两者的取向分布相同——而这个默认在催化体系里通常不成立（表面覆盖度、共吸
附物种、电位都会改变取向）。

\textbf{（3）基底不可复现。}获得活性表面的方法是溶胶化学还原、真空蒸馏与化学沉积，得到的颗
粒尺寸、间隙分布与表面粗糙度难以逐批重现，增强因子因此在批次间漂移。源文本给出的限制更硬：
能产生 SERS 信号的金属和金属氧化物是有限的，研究最多的是 Ag、Au 和 Cu；而且由于在 SERS 研
究中对样品的制备技术要求很高，\textbf{对实际催化剂来说是不容易实现的}，这些使得 SERS 在催化
研究中的进一步应用受到了限制\src{S1 §6.4.4，书内 p.379 / PDF p.398}。

\textbf{要定量只有三条路，且都必须写进图注}：加入内标并测内标峰与目标峰的强度比；在同一种基
底、同一批次上做浓度标定曲线（并给出线性范围）；在同一基底上统计多个点，把热点分布平均掉。
三条都不做，SERS 谱只能回答"有什么物种、以什么取向吸附"，不能回答"有多少"。
\end{pitfall}

\subsubsection{共焦拉曼：剖层分析能剖到多薄}
\label{sec:ch06b-confocal}

共聚焦技术的原理早在 1957 年就已提出，并于 1967 年被用于光学切片分析；1977 年，该技术首次用
于拉曼光谱学\src{S1 §6.4.5，书内 p.379 / PDF p.398}。

\textbf{光路与剖层。}共焦拉曼显微镜的工作原理是：针孔 H1 把激光聚焦于样品表面，样品表面的
被照射点在探测针孔 H2 处成像，其信号由 H2 后的检测器收集；而当激光在样品表面是散焦时，样品
处的大部分信号被 H2 挡住，无法通过针孔到达检测器。把样品沿着激光入射方向上下移动，就可以把
激光聚焦于样品的不同深度，这样所采集的信号也将来自样品的不同深度，实现\textbf{样品的剖层分
析}。这种结构的最大特点就是可以有效地\textbf{排除来自聚焦平面之外其他层信号的干扰}；显微镜
头的数值孔径越大、探测针孔的直径越小，仪器的共焦性能就越好\src{S1 §6.4.5，书内 p.379—380 /
PDF p.398—399}。

\textbf{共焦拉曼的优点}：①高灵敏度，可用于弱拉曼信号的测量；②共聚焦技术的使用使显微镜下激
光在样品上的焦点准确地通过针孔，从而很大程度地提高了\textbf{纵向空间分辨率}，可研究直径为
$1\,$µm 的样品，用于原位多层材料测量；同时共焦显微系统本身还具有较高\textbf{水平方向}的空
间分辨率，而这一分辨率仅取决于显微镜头的放大倍数和所用的激光波长\src{S1 §6.4.5，书内
p.380 / PDF p.399}。

图~\ref{fig:confocal-axial} 用一个模型计算说明"共焦能剖多薄"以及它在什么情况下会失效。

\begin{figure}[tbp]
\centering
\begin{tikzpicture}
\tikzset{declare function={
  nohole(\x) = 1/(1+(\x/8)^2);
  confoc(\x) = 1/(1+(\x/1.5)^2);
  broad(\x)  = 1/(1+(\x/4)^2);
}}
\begin{axis}[
  width=\onecolwidth, height=58mm,
  xmin=-8, xmax=8, ymin=0, ymax=1.05,
  xlabel={离焦量 $z$ / µm},
  ylabel={归一化收集效率},
  xtick={-8,-6,-4,-2,0,2,4,6,8},
  tick label style={font={\fontsize{8}{9.6}\selectfont}},
  label style={font={\fontsize{8}{9.6}\selectfont}},
  legend style={at={(0.02,0.03)}, anchor=south west, draw=none,
     fill=white, fill opacity=0.9, text opacity=1,
     font={\fontsize{8}{9.6}\selectfont}, row sep=0.5pt, inner sep=1.5pt},
  every axis plot/.append style={thick, no markers, samples=250, domain=-8:8},
]
\addplot[black] {confoc(x)};
\addlegendentry{有针孔（理想共焦）}
\addplot[black!60, dashed] {nohole(x)};
\addlegendentry{无针孔（常规显微）}
\addplot[black!30, dotted] {broad(x)};
\addlegendentry{折射率失配或强散射介质}
\end{axis}
\end{tikzpicture}
\caption{模型计算：共焦针孔对轴向（深度）响应的影响\modelmark。读者应看三件事：（i）无针孔时
（虚线）离焦 $\pm5\,$µm 仍收到大部分信号，深度选择性很弱；（ii）加针孔后（实线）只保留焦平
面附近很薄一层的信号，这就是"剖层分析"能成立的物理原因，也是它能测 $1\,$µm 级样品的依据；
（iii）当焦深因样品折射率与物镜标称介质不同、或样品强散射使焦点展宽时（点线），同一针孔的轴
向响应被显著拉宽，深度标度也不再是位移台读数的线性函数——此时"某物种只出现在距表面 $d$ 处"
的结论必须附深度标定数据。三条曲线为洛伦兹型轴向响应的模型计算，半高宽分别取 $16\,$µm、
$3\,$µm 与 $8\,$µm（即参数 $8$、$1.5$、$4$），不代表任何一台仪器或任何一个样品的实测响应；
图例置于左下角并以白底覆盖的仅是三条曲线的尾部（左半支与右半支镜像，形状可由右半支读出）。
\figsrc{剖层原理据 S1 §6.4.5，书内 p.379—380 / PDF p.398—399}}
\label{fig:confocal-axial}
\end{figure}

\textbf{两个应用实例。}

\textbf{（1）混合氧化物催化剂的空间不均一性。}\ce{SEM}-\ce{EDX}（扫描电镜—能量弥散 X 射线）
成像表明 Mo、V 和 W 混合氧化物催化剂中几种元素分布是不均一的；而该催化剂的 XRD 仅给出无定形
相的衍射峰，拟合后表明可能有两种物相存在（无定形 \ce{MoO3} 和无定形 \ce{Mo5O14}）。作者用共
焦拉曼在 $30\times30\,$µm 的催化剂区域内获得了 1000 个拉曼光谱，并用化学计量学方法从 1000 个
光谱中解析出两个主要组分的光谱：组分 1 可能是主要含有 Mo 和 W 的缺氧 \ce{MoWO_{3-x}} 结构，
组分 2 则表明可能是无定形 \ce{V2O5} 的贡献。综合其他分析方法的结果，作者认为在铝基混合氧化
物催化剂中，缺氧的 \ce{MoWO_{3-x}} 组分优先在\textbf{小的高 V 样品区域内}形成，而
\ce{V2O5} 型氧化物是在具有 Mo 和 W 均匀分布的\textbf{剩余区域内}形成的
\src{S1 §6.4.5，书内 p.380—381 / PDF p.399—400}。这个实例的方法论价值在于：\textbf{XRD 给出
的是整块样品的平均物相，共焦拉曼给出的是物相的空间分布}，两者互补而不互相替代——单点 XRD 永
远无法告诉你"哪个组分在哪个微区里"。

\textbf{（2）电极界面的浓度梯度。}田中群等用共焦拉曼显微技术研究了在电催化和燃料电池中有重
要意义的甲醇氧化过程。利用共焦显微拉曼可以剖层分析的特点，他们把激光聚焦于电极表面时，拉曼
信号主要来自电极表面吸附物种；而把激光聚焦在电极表面之上时，拉曼信号则主要来自体相溶液。因
而通过改变激光聚焦点与电极表面的距离，他们得到了反应过程中各物种随离电极表面距离 $d$ 变化的
\textbf{浓度梯度关系}。实验中不仅检测到该体系发生解离吸附的中间产物（CO），还同时监测到溶液
成分的变化——他们的研究表明共焦拉曼光谱可以作为一种多功能的界面研究和分析技术
\src{S1 §6.4.5，书内 p.381—382 / PDF p.400—401}。

\textbf{代价同样存在。}源文本明确指出：由于共焦技术的使用，在一定程度上会破坏样品，并且使原
位检测催化反应变得困难\src{S1 §6.4.5，书内 p.382 / PDF p.401}。这两条限制分别对应高数值孔径
物镜必须在极近距离工作（工作距离短、功率密度高、样品池设计受限），以及共焦光路对窗口与针孔
对准的苛刻要求。

\begin{pitfall}
\textbf{把共焦图上的空间分辨当成物镜的标称分辨。}共焦拉曼的深度选择性由三样东西共同决定，
而它们都不写在"放大倍数"里：
\begin{itemize}
\item \textbf{折射率。}物镜的焦深与横向分辨是按空气（$n\approx1$）标定的；在液相中测量、或
  隔着窗口测样品时，实际焦点位置与标称深度不成正比——介质折射率与物镜设计值不匹配会引入球差，
  使焦点展宽并让深度标度非线性（图~\ref{fig:confocal-axial} 的点线）。用共焦做"浓度随深度
  $d$ 变化"的定量时，\textbf{深度尺度必须用已知标样标定}，不能拿位移台的读数当深度。
\item \textbf{散射体积与多次散射。}针孔只能挡住"离焦"的光，挡不住"在别处被散射之后又拐进针
  孔"的光。催化剂粉末、凝胶与多孔载体都是强散射介质，此时有效采样体积远大于理想共焦体积，横
  向分辨也随之恶化。上文两个实例之所以成立，一个是因为样品是相对致密的氧化物颗粒聚集体，一个
  是因为界面附近的溶液是弱散射介质。
\item \textbf{波长与功率密度。}横向空间分辨率取决于物镜放大倍数与激光波长，紫外激发能给出更
  小的衍射极限，但深紫外物镜的数值孔径与工作距离往往更差；共焦小体积内的功率密度更高，于是又
  回到 \S\,\ref{sec:ch06b-insitu} 的激光加热问题——源文本所说的"共焦技术的使用在一定程度上会
  破坏样品"，正是这一条的实验表现。
\end{itemize}
因此报告共焦结论时必须给出四个参数：\textbf{物镜数值孔径、针孔直径、步长、激发波长与样品介
质}。缺任何一项，别人都无法复现"这个物种出现在哪个深度"。
\end{pitfall}

\subsubsection{应用部分四条误判的汇总}

\begin{pitfall}
本章应用部分最容易犯的四条误判，放在一起对照。它们的共同结构是：\textbf{强度与位置的变化都
存在"非结构"的解释通道}，而排除这些通道只需要改变一个可写进图注的条件。

\begin{itemize}
\item \textbf{把激光加热导致的谱带位移当结构变化。}排除办法是改变激光功率重测并检查时间依赖
  与可逆性（三条判据见 \S\,\ref{sec:ch06b-insitu}）。
\item \textbf{把荧光背底当成宽峰。}宽荧光背景在扣基线之后常留下一个"包"，很容易被当成无定形
  物种或聚合物种的宽峰——分子筛与积炭样品尤其危险，因为模板剂、杂质与碳氢物种都产生极强荧光。
  两条判据：换激发波长重测，真拉曼峰的位移不动、而荧光包的位置随激发波长大幅移动（"峰按位移
  标定、荧光按波长标定"，见本章 §6.2 的"数据到结论"一栏）；测空样品池与纯载体，看这个包是否
  来自样品本身。此外，扣基线方式会直接改变弱峰的相对强度，凡做定量比较必须声明扣基线方法。
\item \textbf{把 SERS 强度直接当浓度。}热点分布、吸附取向与基底批次三者共同决定强度，且能产
  生 SERS 的金属与金属氧化物有限（见 \S\,\ref{sec:ch06b-sers-quant}）。
\item \textbf{把共焦的标称分辨当成实际分辨。}折射率与散射体积都会拉宽轴向响应、改写深度标度
  （见 \S\,\ref{sec:ch06b-confocal}）。
\end{itemize}
凡是没有做这一步排除的谱图，它的归属就只是候选；写成结论时至少要说清"还差哪一步"。
\end{pitfall}

\begin{crosscheck}
\textbf{紫外拉曼、可见拉曼、XRD 与 XPS 给出的不是同一层信息}，这是本节最有价值的一条方法论。
同一块 \ce{ZrO2} 或 \ce{TiO2} 样品，在 \qty{400}{\degC}—\qty{700}{\degC} 的焙烧区间里，XRD
与可见拉曼说体相还是四方相或锐钛矿，紫外拉曼却说表面已经是单斜相、或仍是锐钛矿——两者不矛盾，
因为它们的信息深度不同：紫外强吸收把体相信号衰减掉，剩下的主要是表面区的贡献。

回查规则有四条：
\begin{enumerate}
\item 任何"相变温度"都必须注明是\textbf{体相}还是\textbf{表面}的相变。
\item 当两种技术给出不一致的相变温度时，先不要怀疑哪一方错了，先比较两者的\textbf{信息深度}
  与样品对激发光的吸收系数：紫外区有强吸收的样品，紫外拉曼自然更表面。
\item 需要"表面与体相分开"的结论时，用紫外拉曼（表面）$+$ 可见拉曼或 XRD（体相）成对测量
  ——源文本研究 \ce{TiO2} 相变与光催化活性时用的正是这个组合，最后的活性关联落在\textbf{表面}
  晶相上（表面处于锐钛矿—金红石混合晶相时产氢活性最高，\ce{TiO2(A)/TiO2(R)}-3 的活性约为纯
  金红石载体的 4 倍）。
\item 与本章前半部分对照：那里解决的是"信号从哪来、荧光怎么压"；这里解决的是"信号从多深来、
  能不能比强度"。两组合起来才是完整的判据：\textbf{先确认这是拉曼峰（换波长、扣基线、比空白），
  再确认它来自哪一层（比紫外与可见、标定深度），最后才谈物种与数量。}
\end{enumerate}
把这条规则再往外推一步：本章之外的原位技术（原位 XRD、原位红外、原位电镜）同样各自有信息深
度与原位度，跨技术比较时"条件是否可比"永远优先于"结论是否一致"。
\end{crosscheck}

\paragraph{本章的边界与走向}
从 1928 年拉曼效应的发现算起，拉曼光谱已发展了 70 余年；激光技术的兴起使它成为激光分析中最
活跃的研究领域之一。拉曼光谱在催化研究领域中的优越性在于：可以较容易地实现原位检测，可在高
温高压下表征催化剂的变化，不受气相或载体的干扰。近年来，随着紫外拉曼和共焦显微拉曼光谱等新
的拉曼技术的出现，解决了拉曼光谱早期存在的一些问题，如荧光干扰、固有的灵敏度低等
\src{S1 §6.5，书内 p.402 / PDF p.421}。

限制同样明确，而且都能归到本章的三条线索上：\textbf{荧光与灵敏度}（紫外拉曼与共焦显微拉曼给
出了系统性的解法）、\textbf{激光对样品的改变}（脱水、相变、部分还原甚至完全降解）、以及
\textbf{信息深度}（紫外拉曼因强吸收而表面灵敏，可见拉曼给出体相与表面的混合信息，共焦给出可
标定的深度分辨）。随着纳秒、皮秒激光器的出现及其他光谱技术的不断革新，时间分辨拉曼光谱、近
场共焦拉曼光谱、原子力与拉曼成像联用光谱、光诱导力与拉曼成像联用光谱在催化中的应用成为可能，
这一技术可达到高的时间分辨和空间分辨，从而进一步拓展拉曼光谱在催化中的应用范围，使拉曼光谱
成为传统催化、光催化、材料合成和其他研究领域中一个越来越重要的表征手段
\src{S1 §6.5，书内 p.402—403 / PDF p.421—422}。

对读者的操作性建议只有一条：\textbf{把条件当作结论的一部分来报告}——激发波长、功率密度、光
束尺寸、预处理历史、样品含水量、窗口与介质、针孔直径与步长。本章所有"某物种存在/某物种消失"
的结论，都是在这些条件写全之后才成为证据的。

\begin{methodsum}
\textbf{§6.4—§6.5 的方法骨架}（本分片）。\textbf{应用侧}的读数分三类：骨架与物相（四元环/双
六元环呼吸振动、\ce{MoO3} 与各氧化物晶相特征峰）、负载物种与载体相互作用（低负载量是拉曼最
能发挥作用的区间）、吸附态与孔道内分子（位移与"峰消失"都是信息）。贯穿全章的判据纪律是：
\textbf{谱峰位置只给配位环境，不给唯一物种}。

\textbf{避荧光的两条波长路线}：近红外（FT-Raman，\qty{1064}{nm}）用灵敏度换荧光——$\nu_0^4$
下降、探测器差两个数量级、热辐射干扰、温度上限 $<\qty{150}{\degC}$；深紫外（$<\qty{300}{nm}$）
用选择性与探测深度换荧光与灵敏度——虚态与第一电子激发态之间的间隙把拉曼与荧光在波长上分开，
共振使特定模式增强 $10^2$—$10^6$ 倍，$\nu_0^4$ 再贡献约一个半数量级，而荧光从信噪比的分母
里被删掉。

\textbf{代价侧}：共振使峰高比不再等比；紫外强吸收使探测深度变浅而成为表面灵敏技术（这正是把
表面相与体相分开的工具，\ce{ZrO2} 表面先相变、\ce{TiO2} 表面后相变）；SERS 的强度受热点分布、
吸附取向与基底批次三重控制，只能给物种与取向，不能给浓度；共焦的深度分辨受折射率与散射体积
限制，报结论必须附标定参数；激光加热造成的脱水、相变、还原与降解必须用功率依赖、时间依赖与
可逆性三条判据主动排除。
\end{methodsum}

\srcnote{本节据 S1（上册）第 6 章 §6.4—§6.5，书内 p.373—403 / PDF p.392—422 编写，并收录
§6.3.6 末段一处原位拉曼实例；上一分片已写 §6.1—§6.2（书内 p.343—355 / PDF p.362—374），本节
不再重复拉曼原理、实验技术与荧光发生机制。文中未标注来源的推导与算式（式~\eqref{eq:resonance-cross}
的比例性质说明、波数分辨率与波长分辨率的关系 $\delta\tilde\nu=\delta\lambda/\lambda^2$、以
及图~\ref{fig:uv-fluor-model}、图~\ref{fig:confocal-axial} 两幅模型计算）为本分片作者补写，
非原书内容。源文本为 OCR 文本层，公式、上下标与希腊字母必然失真；本分片对反应式与波数一律采
用"存疑即标注"的处理；定稿时已逐条对照原页面图像核验，原「待核对」标记全部核销，每条
的核验结论见相应位置的\emph{核对说明}。}
